% BHU PYQ -- Manually transcribed & error-corrected
\documentclass[11pt,a4paper]{article}

\usepackage[a4paper,top=2.5cm,bottom=2.5cm,left=2.8cm,right=2.8cm,headheight=14pt]{geometry}
\usepackage{amsmath,amssymb}
\usepackage[T1]{fontenc}
\usepackage[utf8]{inputenc}
\usepackage{lmodern}
\usepackage{microtype}
\usepackage{fancyhdr}
\usepackage{enumitem}
\usepackage{hyperref}
\usepackage{lastpage}
\usepackage{parskip}

\hypersetup{colorlinks=true,urlcolor=black,linkcolor=black,
  pdftitle={BPT-604: Atomic Physics and Lasers -- BHU PYQ},pdfauthor={Banaras Hindu University}}

\pagestyle{fancy}\fancyhf{}
\renewcommand{\headrulewidth}{0.4pt}\renewcommand{\footrulewidth}{0.4pt}
\fancyhead[L]{\small   \textit{Banaras Hindu University}}
\fancyhead[R]{\small   \textit{BPT-604: Atomic Physics and Lasers}}
\fancyfoot[C]{\small Page  \thepage\ of \pageref{LastPage}}


\newlist{parts}{enumerate}{2}
\setlist[parts,1]{label=(\alph*),leftmargin=*,itemsep=5pt,topsep=3pt}
\setlist[parts,2]{label=(\roman*),leftmargin=*,itemsep=3pt,topsep=2pt}

\newcommand{\pts}[1]{\hfill   \textit{[#1]}}


\begin{document}


\begin{center}
  {\large\bfseries BANARAS HINDU UNIVERSITY}\\[2pt]
  {\normalsize B.Sc. (Hons.) Semester VI Examination, 2022-23}\\[6pt]
  \rule{0.8\linewidth}{0.5pt}\\[6pt]
  {\large\bfseries Paper No. BPT-604: Atomic Physics and Lasers}\\[4pt]
  {\normalsize Time: 3 Hours\hfill Maximum Marks: 70}
\end{center}
\rule{\linewidth}{0.4pt}
\smallskip



\noindent   \textit{Instructions: Answer all questions. Symbols have their usual meanings.}
\bigskip




\noindent   \textbf{Question 1.} Answer any   \textbf{five} of the following: \pts{5 $ \times$ 4 = 20}

\begin{parts}
  \item A muonic atom contains a nucleus of one proton and a negative muon $\mu^-$ with charge $-e$ and a mass 207 times as large as an electron mass. What is the wavelength of the first line in the Lyman series for such an atom?
  \item In a Franck-Hertz type of experiment, atomic hydrogen is bombarded with electrons and excitation potentials are found at $10.20\, \text{V}$ and $12.10\, \text{V}$.
    
\begin{parts}
      \item Draw an energy-level diagram and show the possible spectral emissions accompanying these excitations.
      \item What are the wavelengths of the emitted photons?
    \end{parts}
  \item Show that the selection rules for orbital angular momentum and orbital magnetic quantum number are $\Delta l = \pm 1$ and $\Delta m_l = 0, \pm 1$.
  \item A beam of hydrogen atoms emitted from an oven at temperature $400\, \text{K}$ is sent through a Stern-Gerlach magnet of length $1\, \text{m}$. The transverse magnetic field gradient is $10\, \text{T/m}$. Calculate the transverse deflection of the components of the beam at the point where the beam leaves the magnet. Take kinetic energy as $2kT$.
  \item Show the suitability of a two-level atomic system for amplification of radiation in microwave and optical frequency regions.
  \item Evaluate the ratio of the orbital magnetic dipole moment to the orbital angular momentum, $\mu_l/L$, for an electron moving in an elliptical orbit of the Bohr-Sommerfeld atom. Compare the results with those of a circular orbit.
  \item Consider the states in which $l = 4$ and $s = 1/2$. For the state with the largest possible $j$ and largest possible $m_j$, calculate: (i) the angle between $L$ and $S$, (ii) the angle between $\mu_L$ and $\mu_S$, and (iii) the angle between $J$ and the $+z$ axis.
\end{parts}

\medskip\rule{\linewidth}{0.2pt}




\noindent   \textbf{Question 2.}

\begin{parts}
  \item The spin-orbit interaction energy is $\Delta E = \frac{1}{2m^2c^2} \frac{1}{r} \frac{dV(r)}{dr} \vec{S} \cdot \vec{L}$ and the term shift due to relativistic correction of mass for H-atom is $\Delta T_r = \frac{R\alpha^2 Z^4}{n^3} \left[ \frac{1}{l+1/2} - \frac{3}{4n} \right]$. Discuss the fine structure of the $H_\alpha$ line showing clearly the splitting and possible transitions. \pts{8.5}
  \item Compute the ground state spectroscopic terms of carbon and oxygen atoms. How are normal and inverted multiplets different? \pts{4}
\end{parts}

\medskip\rule{\linewidth}{0.2pt}




\noindent   \textbf{Question 3.}

\begin{parts}
  \item What is the Zeeman effect? Derive the expression for interaction energy in a magnetic field much less than the internal magnetic field and apply it to find the components of sodium $D_1$ and $D_2$ lines. \pts{8.5}
  \item Compute the intensity ratio of the three components of the first line of the diffuse series. \pts{4}
\end{parts}

\medskip\rule{\linewidth}{0.2pt}




\noindent   \textbf{Question 4.}

\begin{parts}
  \item Discuss the splitting of $^2P_{3/2}$ and $^2S_{1/2}$ states of sodium atom in a strong magnetic field (greater than the internal magnetic field) and show the transitions. Is it the same as the normal Zeeman effect? \pts{6.5}
  \item In an atom obeying L-S coupling, there are two normal triplets. The separations between adjacent components of the lower triplet are $20\, \text{cm}^{-1}$ and $40\, \text{cm}^{-1}$, and for the higher state triplet the separations are $22\, \text{cm}^{-1}$ and $33\, \text{cm}^{-1}$. Determine the terms for the two triplet states and show the allowed transitions with the help of an energy level diagram. \pts{6}
\end{parts}

\medskip\rule{\linewidth}{0.2pt}




\noindent   \textbf{Question 5.}

\begin{parts}
  \item What are the Einstein coefficients for absorption and emission processes between two levels? Compute the relation between them. \pts{8}
  \item What do you mean by molecular spectra? Explain in brief. \pts{4.5}
\end{parts}
\hfill   \textbf{OR}

\begin{parts}
  \item Discuss the construction and working of a He-Ne Laser. With an energy level diagram, explain the reason for the narrow bore of the Laser tube. \pts{8}
  \item Discuss the rate equations for a three-level Laser system. \pts{4.5}
\end{parts}

\vfill
\rule{\linewidth}{0.4pt}

\begin{center}\small
  \href{https://bhu-exam.vercel.app/}{Click here to visit our website}\\[4pt]
  \href{https://me-aryan.vercel.app/}{Click here to visit our contributor}\\[4pt]
  \href{https://quashan-me.vercel.app/}{Click here to visit our contributor}
\end{center}

\end{document}
